\documentclass[10pt,a4paper]{amsart}
\usepackage[margin=19mm]{geometry}
\usepackage{amsmath,amssymb,mathtools}
\usepackage[scr=rsfs,cal=euler]{mathalfa}
\usepackage{xfrac}
\usepackage[unicode,hidelinks,bookmarks=false]{hyperref}
\newcommand{\ie}{i.e.\ }
\newcommand{\cf}{cf.\ }
\newcommand{\resp}{respectively\ }
\newcommand{\Subsection}{\subsection}
\providecommand{\nobreakdash}{\nobreak\hbox{-}}
\setlength{\emergencystretch}{2em}
\multlinegap0pt
\usepackage{bm}
\usepackage{bbm}
\usepackage{dsfont}
\usepackage{xspace}
\usepackage{cancel}
\usepackage{mathtools}
\usepackage{amsthm}
\makeatletter
\@ifundefined{theorem}{\newtheorem{theorem}{Theorem}}{}
\@ifundefined{lemma}{\newtheorem{lemma}[theorem]{Lemma}}{}
\makeatother
\newcommand{\E}{\mathbb E}
\newcommand{\N}{\mathbb N}
\newcommand{\R}{\mathbb R}
\newcommand{\prob}{\overset{\mathrm{pr.}}{\longrightarrow}}
\newcommand{\weak}{\Longrightarrow}
\renewcommand{\P}{\mathbb P}
\title[The multi-level friendship paradox for sparse random graphs]{The multi-level friendship paradox for sparse random graphs}
\author{Rajat Subhra Hazra, Frank den Hollander, Azadeh Parvaneh}
\date{Source: arXiv:2502.17724v2 (CC BY 4.0)}
\begin{document}
\maketitle

\noindent\emph{Source and licence.} The multi-level friendship paradox for sparse random graphs. Version arXiv:2502.17724v2 (CC BY 4.0), distributed under the Creative Commons Attribution 4.0 International licence, as recorded in the arXiv licence metadata for this exact version and shown on the abstract page https://arxiv.org/abs/2502.17724v2. Full rights notice: licenses/arxiv-2502.17724-CC-BY-4.0.txt.\par\smallskip

\noindent\textit{Editorial scope.} Counted unit: the Lemma whose source label is \texttt{lem} in the article (source0.tex lines 1036--1043, source label \texttt{lem}), delivered together with the article's own complete original proof of it (source0.tex lines 1044--1072, one \texttt{proof} environment of 29 lines). No mathematics is written, repaired or re-derived: statement and proof are the article's own text line for line. Required context retained uncounted: thm1 (lines 493--509); Ginfinity (lines 854--862); hindicator (lines 901--904). Every definition, display and label of those lines is retained unchanged. Omitted from this package and not redistributed: 1--492, 510--853, 863--900, 905--1035, 1073--1540. Nothing inside a retained excerpt is omitted or altered: every retained body is delivered exactly as the article's lines are written, so no omission receipt of any kind accompanies this package. Citations: the retained excerpts carry no citation command at all, so no bibliography entry of the article is transcribed and no reference is redistributed. Reference adaptation: none was needed and none was made; every reference inside the delivered unit resolves inside it, so no unresolved reference or citation is produced. Printed slips kept literal: no printed word, symbol, hypothesis or number of the counted statement or of its proof was altered. Layout and class: the article's own class is \texttt{article} with packages including babel, inputenc, amsfonts, amssymb, amsmath, amsthm, latexsym, epsfig, amscd, bbm, stmaryrd, mathrsfs, hyperref, fancyhdr; the wrapper is the lane's pinned \texttt{amsart} editorial wrapper at 10pt on A4, which restates the theorem-like environments the retained text needs and adds the article's own macro definitions that the retained text needs (\texttt{\textbackslash E}, \texttt{\textbackslash N}, \texttt{\textbackslash P}, \texttt{\textbackslash R}, \texttt{\textbackslash prob}, \texttt{\textbackslash weak}), with extra wrapper packages (bm, bbm, dsfont, xspace, cancel, mathtools). The delivered numbering is the unit's own. Assets: no retained span contains a figure environment, a graphics-inclusion command, a PDF-page inclusion, a table or a TikZ picture; no image file is copied into this package, the package declares no asset dependency and no image file is redistributed with it. Licence: version arXiv:2502.17724v2 of this article is distributed under the Creative Commons Attribution 4.0 International licence, as recorded in the arXiv licence metadata for this exact version and shown on the abstract page https://arxiv.org/abs/2502.17724v2. A full rights notice is delivered at licenses/arxiv-2502.17724-CC-BY-4.0.txt. Characters: every retained excerpt is pure ASCII, so the delivered engine, XeLaTeX with its default Latin Modern text font, reports no missing character.
\begin{theorem}
\label{thm1}
For both types of exploration and each $k\in\N$:
\begin{itemize}
\item[\rm{(a)}] 
If $G_{n}$ converges locally in probability to $(G_{\infty},\phi)$ as $n\to\infty$, then $\mu_{n}^{(k)} \weak \mu^{(k)}_{\infty}$ as $n\to\infty$ in probability. In particular, as $n\to\infty$,
\begin{align*}
\mu_{n}^{(k)} (A) \prob \mu^{(k)}_{\infty} (A), \qquad 
\tilde{\mu}_{n}^{(k)} (A) \to \mu^{(k)}_{\infty} (A) \qquad \forall\,A \in \mathcal{B}(\R).
\end{align*}
\item[\rm{(b)}] 
If $G_{n}$ converges locally weakly to $(G_{\infty},\phi)$ as $n\to\infty$, then $\tilde{\mu}_{n}^{(k)} \weak \mu^{(k)}_{\infty}$ as $n\to\infty$. In particular, as $n\to\infty$,
\begin{align*}
\tilde{\mu}_{n}^{(k)} (A) \to \mu^{(k)}_{\infty} (A) \qquad \forall\,A \in \mathcal{B}(\R).
\end{align*}
\end{itemize}
\end{theorem}

\begin{align}
\label{Ginfinity}
(G_{\infty}^{\prime}(\omega),\phi) = 
\left\{
\begin{array}{ll}
(G_{\infty}(\omega),\phi), &\text{if }(G_{\infty}(\omega),\phi) \text{ is locally finite}, \\
(H,\phi), &\text{otherwise},
\end{array} \right.
\end{align}

\begin{align}
\label{hindicator}
h((G,o))=\mathbbm{1}_{\{\Delta_{o,k}^{(G)}\in A\}},
\end{align}

\begin{lemma}
\label{lem}
If $(H_{n})_{n\in\N}$ is a sequence of finite random graphs with a deterministic vertex set that converges locally in probability to an almost surely locally finite, connected rooted graph $(H, \varphi)$, and if $((d_{i}^{(H_{n})})^{m})_{i\in V(H_{n}),\, n\in\N}$ is uniformly integrable for some $m \in \N$, then
\begin{align*}
\frac{1}{\#V(H_{n})}\sum_{i\in V(H_{n})}(d_{i}^{(H_{n})})^{j}\prob \E_{H}\big[(d_{\varphi}^{(H)})^{j}\big],\qquad n\to\infty ,
\end{align*}
for all $j\in\{1,\ldots,m\}$, where $\E_{H}$ denotes the expectation with respect to the randomness of the graph $(H, \varphi)$. 
\end{lemma}

\begin{proof}
Consider $j \in \{1, \ldots, m\}$. Without loss of generality, assume that $(H, \varphi)$ is everywhere locally finite, in a manner similar to the construction in \eqref{Ginfinity}. 

For $M \in \N$, define a bounded function $h_{j,M}\colon\,(\mathscr{G},\mathrm{d}_{\mathscr{G}}) \to (\R,\vert\cdot\vert)$ by setting
\begin{align*}
h_{j,M}((G,o))=(d^{(G)}_{o})^{j}\,\mathbbm{1}_{\{d^{(G)}_{o}\leq M\}}.
\end{align*}
Then, in a manner analogous to the proof of the continuity of the function in \eqref{hindicator} provided in the proof of Theorem \ref{thm1}, it can be easily shown that $h_{j,M}$ is continuous. Therefore, by the local convergence in probability of $(H_{n})_{n\in\N}$ to $(H, \varphi)$, we have
\begin{align}\label{lem1}
\frac{1}{\#V(H_{n})}\sum_{i\in V(H_{n})}(d_{i}^{(H_{n})})^{j}\,\mathbbm{1}_{\{d_{i}^{(H_{n})}\leq M\}}\prob \E_{H}\Big[(d_{\varphi}^{(H)})^{j}\,\mathbbm{1}_{\{d_{\varphi}^{(H)}\leq M\}}\Big],\qquad n\to\infty .
\end{align}
Moreover, by the monotone convergence theorem, we have that
\begin{align*}
\E_{H}\Big[(d_{\varphi}^{(H)})^{j}\,\mathbbm{1}_{\{d_{\varphi}^{(H)}\leq M\}}\Big]\to\E_{H}\big[(d_{\varphi}^{(H)})^{j}\big]\qquad M\to\infty .
\end{align*}
Therefore, for any $\epsilon >0$, we can take $M_{\epsilon}\in\N$ such that for all $M\geq M_{\epsilon}$,
\begin{align*}
\bigg|\E_{H}\Big[(d_{\varphi}^{(H)})^{j}\,\mathbbm{1}_{\{d_{\varphi}^{(H)}\leq M\}}\Big]-\E_{H}\big[(d_{\varphi}^{(H)})^{j}\big]\bigg|<\frac{\epsilon}{4},
\end{align*}
and by the Markov inequality,
\begin{align*}
&\P_{H_{n}}\bigg\{\Big |\tfrac{1}{\#V(H_{n})}\sum_{i\in V(H_{n})}(d_{i}^{(H_{n})})^{j}-\E_{H}\big[(d_{\varphi}^{(H)})^{j}\big]\Big |>\epsilon\bigg\}
\\&\leq \P_{H_{n}}\bigg\{\Big |\tfrac{1}{\#V(H_{n})}\sum_{i\in V(H_{n})}(d_{i}^{(H_{n})})^{j}\,\mathbbm{1}_{\{d_{i}^{(H_{n})}\leq M\}}-\E_{H}\big[(d_{\varphi}^{(H)})^{j}\big]\Big|>\tfrac{\epsilon}{2}\bigg\}\\&\quad
+\P_{H_{n}}\bigg\{\tfrac{1}{\#V(H_{n})}\sum_{i\in V(H_{n})}(d_{i}^{(H_{n})})^{j}\,\mathbbm{1}_{\{d_{i}^{(H_{n})}> M\}}>\tfrac{\epsilon}{2}\bigg\}\\&
\leq \P_{H_{n}}\bigg\{\Big |\tfrac{1}{\#V(H_{n})}\sum_{i\in V(H_{n})}(d_{i}^{(H_{n})})^{j}\,\mathbbm{1}_{\{d_{i}^{(H_{n})}\leq M\}}-\E_{H}\big[(d_{\varphi}^{(H)})^{j}\,\mathbbm{1}_{\{d_{\varphi}^{(H)}\leq M\}}\big]\Big|>\tfrac{\epsilon}{4}\bigg\}\\&\quad
+\frac{2}{\epsilon}\,\sup_{n\in\N}\,\sup_{i\in V(H_{n})}\E_{H_{n}}\bigg[(d_{i}^{(H_{n})})^{m}\,\mathbbm{1}_{\{d_{i}^{(H_{n})}> M\}}\bigg],
\end{align*}
where $\P_{H_{n}}$ denotes the probability measure on the underlying space of the random graph $H_{n}$, and $\E_{H_{n}}$ represents the expectation with respect to the randomness of $H_{n}$. By taking the limit as $n\to\infty$ and then $M\to\infty$, and applying \eqref{lem1} along with the uniform integrability condition, we establish the lemma.
\end{proof}

\end{document}
